\section{Discussion}\label{discussion}

The results show that OPSD did not consistently outperform conventional SFT for visual question
answering. For SmolVLM and BLIP, SFT achieved substantially higher BertScore-F1 and exact-match
performance, while MiniCPM showed similar results between the two approaches which can be seen in \autoref{tab:performance}. These findings
provide a negative answer to RQ1 under the experimental conditions considered: privileged-context
OPSD alone was not generally more effective than direct supervised fine-tuning. However, the magnitude
of the difference varied considerably across architectures, providing a more nuanced answer for RQ2.

\subsection{Why SFT Outperformed OPSD}

The strongest explanation for the performance gap is the difference in supervision available to the
two approaches. SFT receives direct supervision from the ground-truth answer at every training
step. The optimisation objective therefore explicitly increases the likelihood of the correct target
sequence. In contrast, OPSD condition used in this study receives no direct ground-truth supervision. Instead,
the student is trained only to match the token-level distribution produced by the privileged
teacher on the student's own generated trajectory.

This observation becomes particularly important when the privileged teacher is only marginally better
than the student. Prior to fine-tuning, the crop-conditioned teacher advantage was measured for each
architecture. The BertScore-F1 advantage of the crop over the full image was $-0.00067$ for SmolVLM,
$+0.00188$ for BLIP and $-0.00027$ for MiniCPM. These differences indicate that the privileged
crop did not produce a consistently stronger teacher across architectures. Although the crop sometimes
improved exact-match accuracy, its average advantage was small.

OPSD relies on useful information being present in the difference between the teacher and student
distributions. If the privileged teacher behaves similarly to or on occasion worse than the full-image
student, this can result in the distillation signal providing limited information beyond the model's 
existing behaviour. SFT, on the other hand, does not face this limitation since its supervisory signal
is obtained from the labelled target (ground-truth answer). The relatively weak teacher advantage
therefore provides a plausible explanation as to why OPSD remained close to the pretrained baseline for
SmolVLM and BLIP while SFT produced better results.

\subsection{Architecture-Dependent Behaviour}

Although SFT was stronger overall, the relative OPSD-SFT gap decreased substantially across the
three architectures. From \autoref{tab:paired_differences}, it can be seen that the paired BERTScore-F1 difference improved from (-0.0453) for 
SmolVLM to (-0.0157) for BLIP and finally to approximately zero ((-0.0001)) for 
MiniCPM. For MiniCPM, the 95\% bootstrap confidence interval crossed zero and OPSD 
achieved slightly higher exact-match accuracy than SFT.

The per-example analysis seen in \autoref{tab:per_example_comparison} reinforces this difference.
SmolVLM showed a strong preference towards SFT, which achieved a higher BERTScore-F1 on 
81.77\% of examples. In contrast, MiniCPM produced identical scores for 82.68\% of test 
examples, with OPSD and SFT winning on 8.33\% and 8.98\% respectively. This suggests 
that OPSD and SFT converged to substantially more similar behaviour for MiniCPM than 
for the smaller models.

These results suggest that the effectiveness of privileged-context self-distillation may depend
on model characteristics. A stronger multimodal model such as MiniCPM, may already possess
better visual grounding and reasoning capabilities, allowing it to make more effective use of
subtle distributional differences between full-image and crop-conditioned contexts. On the other hand,
a smaller model may benefit more from explicit labelled supervision than from a comparatively weak
teacher signal. Importantly since only three architectures were evaluated, this pattern should not be used
to infer a general relationship between model size and OPSD effectiveness but rather than indicator
that OPSD performance cannot be assumed to transfer uniformly across VLM architectures.

\subsection{Local Evidence Versus Global Context}

An important limitation of the current OPSD setting is that the privileged visual context
is only as useful as the crop from which it is constructed. Although GQA provides grounded
object annotations, these annotations were not originally designed specifically for privileged-context 
self-distillation. As a result, the generated crop may contain the target object without necessarily
containing all of the visual evidence required to answer a question. The relatively weak teacher
advantage observed before training supports this limitation. \autoref{fig:qualitative_examples} illustrates
two images that represent this trade-off:

\begin{figure}[h!]
    \centering
    
    \begin{subfigure}{0.48\columnwidth}
        \centering
        \includegraphics[width=\linewidth]{images/opsd_example.png}
        \caption{OPSD-favouring example.}
        \label{fig:opsd}
    \end{subfigure}
    \hfill
    \begin{subfigure}{0.48\columnwidth}
        \centering
        \includegraphics[width=\linewidth]{images/sft_example.png}
        \caption{SFT-favouring example.}
        \label{fig:sft}
    \end{subfigure}

    \caption{Representative MiniCPM examples illustrating the trade-off between localised privileged evidence and global scene context.}
    \label{fig:qualitative_examples}
\end{figure}

In \autoref{fig:opsd}, OPSD correctly answers a question regarding the colour of the kite-surfer's 
attire while the SFT produces the incorrect response. In this situation, the relevant visual evidence
occupies only a relatively small portion of a rather visually complex image. In this case, the
crop-conditioned teacher provides a more focused view of the relevant evidence, which is consistent with the
intended benefit of privileged visual context. In contrast, \autoref{fig:sft} represents a case
in which SFT correctly answer a question asking how large the boat was which OPSD got incorrect. Unlike
the previous example, judging relative size requires information from the wider scene. By restricting
the teacher to a local crop may therefore remove the contextual references required to make the correct
judgement, even when padding is included.

These two examples demonstrate that successful crop-based OPSD requires more than accurately localising
the target object. The privileged context must preserve the complete evidence needed to answer the question.
A crop than therefore be beneficial when the task depends mainly on fine-grained local evidence, but 
potentially harmful when the required reasoning depends on global or relational context.

\subsection{Limitations and Future Work}

Several limitations should be considered. First, the experiment used a 5000-examle subset of
GQA, which may disadvantage OPSD because it relies on indirect distributional supervision rather than
direct target labels. In addition, additional datasets, architectures and random seeds 
would also be required before drawing broader conclusions about OPSD performance across VLMs.

Second, the privileged teacher used deterministic question-grounded crops. Although unsuitable
spatial and relational questions were filtered, GQA was not designed specifically for crop-based privileged
supervision. This implies that some crops may omit context required to answer a question. Future work
could investigate multiple evidence regions, learned region selection, or combined local and global teacher views.

Third, the OPSD condition used no supervised cross-entropy anchor. A useful extension would therefore be to
combine OPSD with a small supervised objective and evaluate whether this improves stability while
retaining the benefit of privileged information.

Lastly, OPSD may be more effective in domains where the privileged context provides a stronger
information advantage. For example, in medical imaging, a teacher could receive a radiologist-defined
region of interest or lesion annotation while the student operated on the complete image. 